\begin{itshape}
The majority of websites continue to run PHP web applications on end-of-life versions that no longer receive official security updates, leaving known vulnerabilities permanently unpatched. Manual migration of hundreds of source files is impractical, and existing transformation tools do not incorporate security compliance checks as an integrated part of the migration workflow. Organizations operating under ISO/IEC 27001:2022 are further burdened by having to conduct security audits as a separate step after migration is complete.

\hspace{1cm} This research develops a Python-based platform that integrates automated conversion of legacy PHP web applications to a configurable target PHP version with ISO/IEC 27001:2022-aligned secure coding compliance checks in a single, fully local, open source workflow. The platform combines Rector for AST-based code transformation, Semgrep for security scanning, PHPStan for post-conversion type validation, and five locally executed open source large language models via Ollama---DeepSeek Coder 6.7b, Qwen2.5-Coder 7b, CodeLlama 7b, Mistral 7b, and Llama 3.1 8b---for semantic vulnerability analysis and remediation suggestions. Case studies were conducted on PHP codebases from FT UGM (Dashboard TCK, 245 files) and CBT DTI UGM (326 files).

\hspace{1cm} Rector achieved a conversion rate of 97.6\% on the FT UGM dataset and 91.7\% on the CBT DTI UGM dataset, with 100\% post-conversion syntax validity in both cases. A delta finding value of zero ($\Delta F = 0$) confirmed that the conversion process introduced no new security vulnerabilities. Automatically generated JSON reports include per-finding mappings to relevant ISO/IEC 27001:2022 Annex A controls and can be used directly as audit documentation. Among the five evaluated LLMs, Qwen2.5-Coder 7b is recommended for automated pipeline integration, consistently achieving 100\% Syntax Validity Rate across all experimental conditions. The platform demonstrates that automated PHP version migration combined with ISO/IEC 27001:2022 compliance reporting can be executed entirely on-premises without transmitting source code to any external service.
\end{itshape}

\noindent\textbf{Keywords} : PHP migration, secure coding, ISO/IEC 27001:2022, large language model, AST code transformation
